\chapter{Our Greedy Algorithms}
\label{ch:our-algorithms}
In this section, we introduce our methods for synthesizing low-depth CNOT circuits for block ciphers. Our approach is based on Shi and Feng's greedy method, specifically for block sizes of $16 \times 16$ and $32 \times 32$, as discussed in \cite{cryptoeprint:2023/165}.

We propose a parameterized, depth-oriented algorithm for synthesizing ancilla-free, low-depth CNOT circuits, and we develop two distinct greedy algorithms: Row-or-Column and Parallel, to further explore depth-optimized results. By utilizing different matrix representations and heuristic search strategies, the implementation logic and pseudocode for each approach are detailed in the following sections.

\section{Proposed Approach}

The primary idea integrating these cost metrics and greedy heuristics 
is detailed in Algorithm~\ref{alg:core-idea}. To elaborate on our approach, 
we first initialize the upper bounds ``minm`` and ``minm\_size`` 
for the subsequent calculations. At the start of the process, a copy of 
the input matrix is retained for final verification. Simultaneously, the 
matrix inverse is calculated to facilitate evaluation of the cost function. 
In addition, we initialize all required lists, and ``over`` is set to 
\textbf{False}. After all variables have been defined, we execute the 
greedy algorithm, using the value of $p$ to select the operation type and 
cost function.

Once the greedy algorithm returns the required operations, if ``over`` 
is triggered, the current iteration is terminated immediately and the next 
iteration is prepared; otherwise, we proceed to arrange the available 
operators. First, we check the operator lists $L_{r}$ and $L_{c}$ from the 
previous iteration, as some operations may remain that do not fill the 
list but can still reduce the cost. In this case, we append them to 
$Ls_{r}$ and $Ls_{c}$, as shown in Algorithm~\ref{alg:last-check}, and 
subsequently clear $L_{r}$ and $L_{c}$.

Next, to guarantee that all collected operators are valid and to count 
their number for the circuit, as shown in Algorithm~\ref{alg:verify-operations}, 
we use ``Verify\_Matrix()`` to copy $M$ together with $OPs_{r}$ and 
$OPs_{c}$. This procedure applies ``row\_i2j()`` and ``col\_i2j()`` 
to reduce the replica while simultaneously recording the number of 
operations used. If the reduced result is not equal to $M$, this indicates that the
sequence of CNOT gates is invalid; otherwise, if the initial matrix is
transformed into $M$ after the CNOT gates have been applied, the
operators are valid, and the corresponding CNOT gates are identified as
effective.

We then check whether the depth $d$ exceeds the maximum value, or whether 
the CNOT count exceeds the maximum size. If either condition holds, the 
operators are considered invalid, and we break out of the current 
iteration. Otherwise, the depth and CNOT count are valid, and we update 
the current record accordingly.

Once all required operators have been obtained, we record them to local 
files, as shown in Algorithm~\ref{alg:recorder}. We first create the 
permutation list and then record the sequence of CNOT gates. Specifically, 
the order of operations recorded during the reduction of $M$ to a 
permutation matrix $P$ must be reversed when synthesizing the actual CNOT 
circuit for $M$. This follows directly from matrix inversion: since 
$M = (E^d_{i_d}\cdots E^1_1)^{-1} \cdot P \cdot (F^1_1\cdots F^d_{j_d})^{-1}$, 
and inverting a product reverses the order of its factors, the elementary 
operations used to reduce $M$ to $P$ must be applied in the opposite order 
to reconstruct $M$ from $P$. Because each type-3 elementary matrix over 
$\mathbb{F}_2$ is self-inverse, this reversal requires no additional matrix 
inversion, only a reordering of the recorded operation sequence. Accordingly, 
``sequence`` and ``Layers`` are constructed following this 
principle. Overall, we generate files containing the permutation list, the 
CNOT gate operations, and each layer for the input matrix.

Finally, to globally verify that ``sequence`` and ``Layers`` are 
correct, we develop the ``Verify\_All()`` function, as shown in 
Algorithm~\ref{alg:verify-all}. The verification begins by checking that 
each individual layer in ``Layers`` is free of conflicts, ensuring 
that all operators within the same layer are distinct. Once we have 
confirmed that no operations are duplicated, we execute them on a copy of 
$M$ to verify that they indeed reproduce the final state of $M$. We further 
ensure that all operators in ``Layers`` are unique by maintaining a 
layer-visited list. We then verify that ``sequence`` correctly applies 
its CNOT operations to ``origin``. Rather than replaying the recorded 
sequence in reverse order to reconstruct $M$ directly, verification instead 
exploits the duality between row and column operations established in 
Section~\ref{sec:CNOT-transformation-to-Matrix-Operations}: a row operation 
$(control, target)$ applied to $M$ is equivalent to a column operation 
$(target, control)$ applied to $M^{-1}$, with the roles of control and 
target exchanged. This allows the recorded sequence to be replayed in its 
original, forward order, applying each row-type entry as a column operation 
on the dual matrix with its control and target swapped, and each 
column-type entry as a row operation with the same swap. This avoids the 
need to explicitly reverse the operation list during verification. Through 
this verification procedure, if ``origin`` is transformed into $M$, 
we can confirm that all operations are valid for execution on a quantum 
simulator, and the algorithm returns the depth $d$ and ``Layers``.

\begin{algorithm}[htbp]
\caption{Core Idea}
\label{alg:core-idea}
\begin{algorithmic}[1]
\Require $M \in \mathbb{F}_{2}^{n \times n}$ and $p$
\Ensure Calculate the depth $d$ of the matrix and record all available operators in Layers
\State minm, minm\_size $\leftarrow$ MAX \Comment{Assign the maximum value supported by the operating system as the initial upper bound.}
\State $inv \leftarrow$ inverse($M$) \Comment{Calculate the initial inverse for $M$.}
\State $L_{r}, L_{c}, Ls_{r}, Ls_{c}, OPs_{r}, OPs_{c} \leftarrow$ list[$\varnothing$] \Comment{Define all used lists.}
\State over $\leftarrow False$ \Comment{Initialize the depth-based break condition.}

\State $L_{r}, L_{c}, Ls_{r}, Ls_{c}, M, OPs_{r}, OPs_{c}, d$, over $\leftarrow$ Row/Col/Row-or-Col/Parallel Greedy($p$) \Comment{Select the greedy algorithm and cost function according to $p$.}

\If{over}
	\State \Return \Comment{The greedy algorithm checks whether the depth of $M$ is greater than the latest minimum depth.}
\EndIf

\State $Ls_{r}, Ls_{c} \leftarrow$ Last\_Check($L_{r}, L_{c}$) \Comment{Check the last collection operations.}

\State $size, ok$ $\leftarrow$ Verify\_Matrix($OPs_{r}, OPs_{c}, M$) \Comment{Check whether the operations work on $M$.}

\If{($d > minm$) || ($d == minm \land size \geq minm\_size$) || $!ok$}
	\State \Return \Comment{Check whether the depth and CNOT count exceed the current minimum recorded values.}
\EndIf

\State $minm \leftarrow d$ \Comment{Update depth $d$.}
\State CNOT\_count $\leftarrow size$ \Comment{Update CNOT count.}

\State per, seq, Layers $\leftarrow$ Recorder($M, OPs_{r}, OPs_{c}, Ls_{r}, Ls_{c}$) \Comment{Create permutation, CNOT sequence and operations Layer.}

\State correct $\leftarrow$ Verify\_All($M$, Layers, seq) \Comment{Verify whether the CNOT sequence and operation layers can transform the origin into $M$.}
\If{!correct}
	\State \Return \Comment{Break if the operations do not work for the origin to become $M$.}
\EndIf
\State \Return $d$, Layers
\end{algorithmic}
\end{algorithm}

\begin{algorithm}[htbp]
\caption{Last\_Check}
\label{alg:last-check}
\begin{algorithmic}[1]
\Require $L_{r}, L_{c}$
\Ensure Find the last round operations and insert them to $Ls_{r}$ and $Ls_{c}$
\If{$|L_{r}| > 0$}
	\State $Ls_{r}$.append($L_{r}$) \Comment{Insert the remaining $L_{r}$ into $Ls_{r}$.}
	\State $L_{r}$.clear() \Comment{Clear $L_{r}$.}
\EndIf
\If{$|L_{c}| > 0$}
	\State $Ls_{c}$.append($L_{c}$) \Comment{Insert the remaining $L_{c}$ into $Ls_{c}$.}
	\State $L_{c}$.clear() \Comment{Clear $L_{c}$.}
\EndIf
\State \Return $Ls_{r}, Ls_{c}$
\end{algorithmic}
\end{algorithm}

\begin{algorithm}[htbp]
\caption{Verify\_Matrix}
\label{alg:verify-operations}
\begin{algorithmic}[1]
\Require $M_{origin} \in \mathbb{F}_{2}^{n \times n}, M \in \mathbb{F}_{2}^{n \times n}, OPs_{r}, OPs_{c}$
\Ensure Check operations from $OPs_{r}$ and $OPs_{c}$ that can work on $M$
\State reduce $\leftarrow$ $M_{origin}$ \Comment{Copy the initial matrix $M_{origin}$ without applying any CNOT gates.}
\State size $\leftarrow 0$ \Comment{Initialize the CNOT count.}
\For{$i = 0$ \textbf{to} $|OPs_{r}|-1$}
	\State reduce $\leftarrow$ row\_i2j(reduce, $OPs_{r}[i]$.control, $OPs_{r}[i]$.target) \Comment{Apply the row operation to ``reduce``.}
	\State size++ \Comment{Increase the CNOT gate count.}
\EndFor 
\For{$j = 0$ \textbf{to} $|OPs_{c}|-1$}
	\State reduce $\leftarrow$ col\_i2j(reduce, $OPs_{c}[j]$.control, $OPs_{col}[j]$.target) \Comment{Apply the column operation to ``reduce``.}
	\State size++ \Comment{Increase the CNOT gate count.}
\EndFor
\State ok $\leftarrow True$
\If{reduce $\neq M$} \Comment{Verify that ``reduce`` equals $M$ after applying the CNOT gates.}
	\State ok $\leftarrow False$ \Comment{Break if ``reduce`` is not equal to $M$.}
\EndIf 
\State \Return size, ok
\end{algorithmic}
\end{algorithm}

\begin{algorithm}[htbp]
\caption{Recorder}
\label{alg:recorder}
\begin{algorithmic}[1]
\Require $M \in \mathbb{F}_{2}^{n \times n}, OPs_{r}, OPs_{c}, Ls_{r}, Ls_{c}$
\Ensure Collection permutation list, sequence list and Layers 
\State per, seq, Layers $\leftarrow [\varnothing]$ \Comment{Define the required lists.}
\For{$i = 0$ \textbf{to} $|M|-1$}
	\For{$j = 0$ \textbf{to} $|M|-1$}
		\If{$M[i][j]$}
			\State per[$i$] = $j$ \Comment{Collect the permutation list.}
		\EndIf
	\EndFor
\EndFor 
\For{$i = 0$ \textbf{to} $|OPs_{c}|-1$}
	\State seq.append($OPs_{c}[i]$.target, $OPs_{c}[i]$.control) \Comment{Collect the available operations into a sequence.}
\EndFor
\State $OPs_{r} \leftarrow$ reverse($OPs_{r}$) \Comment{Reverse $OPs_{r}$ with respect to the column operations.}
\For{$j = 0$\ \textbf{to} $|OPs_{r}|-1$}
	\State seq.append(per[$OPs_{r}[j]$.control], per[$OPs_{r}[j]$.target]) \Comment{Record the row operations.}
\EndFor
\For{$i = 0$ \textbf{to} $|Ls_{c}|-1$}
	\State nl\_c $\leftarrow [\varnothing]$ \Comment{Create a single column layer for ``Layers``.}
	\For{$j = 0$ \textbf{to} $|Ls_{c}[i]|-1$}
		\State nl\_c.append($Ls_{c}[i][j]$.target, $Ls_{c}[i][j]$.control) \Comment{Record a single column layer.}
	\EndFor
	\State Layers.append(nl\_c) \Comment{Append a single column layer to ``Layers``.}
\EndFor
\State $Ls_{r} \leftarrow$ reverse($Ls_{r}$) \Comment{Reverse $Ls_{r}$ with respect to the column layer.}
\For{$k = 0$ \textbf{to} $|Ls_{r}|-1$}
	\State nl\_r $\leftarrow [\varnothing]$ \Comment{Create a single row layer for ``Layers``.}
	\For{$l = 0$ \textbf{to} $|Ls_{r}[k]|-1$}
		\State nl\_r.append(per[$Ls_{r}[k][l]$.control], per[$Ls_{r}[k][l]$.target]) \Comment{Record a single row layer.}
	\EndFor
	\State Layers.append(nl\_r) \Comment{Append a single row layer to ``Layers``.}
\EndFor
\State \Return per, seq, Layers
\end{algorithmic}
\end{algorithm}

\begin{algorithm}[htbp]
\caption{Verify\_All}
\label{alg:verify-all}
\begin{algorithmic}[1]
\Require $M_{origin} \in \mathbb{F}_{2}^{n \times n}, M \in \mathbb{F}_{2}^{n \times n}$, Layers, seq
\Ensure Check all operations are work
\State origin $\leftarrow M_{origin}$ \Comment{Copy the initial matrix without applying any CNOT gates.}
\If{!Conflict\_Check(Layers)} \Comment{Check that all layers in ``Layers`` are distinct.}
	\State \Return $False$
\EndIf
\For{$i = 0$ \textbf{to} |Layers|-1}
	\State layer\_visi $\leftarrow [\varnothing]$ \Comment{Create a layer-visited list.}
	\For{$j = 0$ \textbf{to} |Layers[i]|-1}
		\If{layer\_visi[Layers[i][j].control] or layer\_visi[Layers[i][j].target]}
			\State \Return $False$ \Comment{Check whether the CNOT control and target exist.}
		\Else
			\State layer\_visi[Layers[i][j].control] $\leftarrow 1$ \Comment{Register the control in the layer-visited list.}
			\State layer\_visi[Layers[i][j].target] $\leftarrow 1$ \Comment{Register the target in the layer-visited list.}
		\EndIf
	\EndFor
\EndFor
\For{$k = 0$ \textbf{to} |seq|-1}
	\If{seq[k].op == row}
		\State origin $\leftarrow$ col\_i2j(origin, seq[k].target, seq[k].control) \Comment{Apply the reversed sequence of row operations to ``origin``.}
	\Else
		\State origin $\leftarrow$ col\_i2j(origin, seq[k].control, seq[k].target) \Comment{Apply the column operations from ``sequence`` to ``origin``.}
	\EndIf
\EndFor
\If{origin == $M$} \Comment{Check whether the origin matrix is the same as the CNOT-applied matrix $M$.}
	\State \Return $True$
\Else
	\State \Return $False$
\EndIf
\end{algorithmic}
\end{algorithm}

\subsection{Row-or-Col Greedy Algorithm}
The first strategy evaluates both row and column operations and selects 
the layer that yields the minimum matrix cost, as shown in 
Algorithm~\ref{alg:row-or-col-greedy}. In the spirit of the aforementioned 
row/column greedy algorithm, our goal remains to collect available 
operators for $Ls_{r}$ and $Ls_{c}$; however, in this case, we explore 
possible operations bidirectionally. Specifically, we use 
$best_{r}^{cost}$ and $best_{c}^{cost}$, which maintain a collection of 
cost values for each potential operation in $B_{r}$ and $B_{c}$, to 
facilitate selection of the optimal operator.

To address the inherent risk of local minima in greedy heuristics, we 
maintain a collection of sub-optimal operators, representing candidates 
whose cost values exceed the current minimum. A counter is used to 
record the number of consecutive stalls (i.e., stuck iterations) in 
order to detect when the algorithm has reached a local optimum. A 
stagnation threshold is defined to limit the number of consecutive 
stalls; by default, this value is set to three. If the cost fails to 
decrease for three consecutive iterations, the algorithm accepts a 
sub-optimal operation that increases the cost, thereby enabling an 
escape from the local minimum. That is, if neither the row nor column 
operators can achieve a cost decrease within the threshold (three 
iterations by default), the stall counter is incremented accordingly. 
In this case, we execute the ``Rand\_Ops\_Sel()`` function, which 
is similar to Algorithm~\ref{alg:ops-sel} but without the ``minm`` 
constraint. We then identify the minimum-cost candidate returned by 
``Rand\_Ops\_Sel()``, insert it into ``select\_list``, and 
update ``minm`` accordingly for this cost-increasing case.

After the required operators have been collected, we compare 
$best_{r}^{cost}$ and $best_{c}^{cost}$ to determine which operator list 
should be selected, that is, which layer achieves the greater cost 
reduction. In cases where the two costs are identical, the algorithm 
defaults to the column-wise operation. If only one set of operations is 
available, the algorithm selects that existing set.

We then execute the operations from ``select\_list``. Unlike the 
row/column greedy algorithm, in the row-or-column strategy we switch the 
cost function to $h_{prod}(M)$ once ``close\_permu`` is triggered, in 
order to handle sparse matrices; this achieved a depth of 9 and 109 gates 
for AES MixColumns. Finally, we check whether the depth exceeds the 
previous minimum cost and set the ``over`` flag accordingly.

\begin{algorithm}[htbp]
\caption{Row-or-Col Greedy}
\label{alg:row-or-col-greedy}
\scriptsize
\begin{algorithmic}[1]
\Require $M \in \mathbb{F}_{2}^{n \times n}$, $M' \in \mathbb{F}_{2}^{n \times n}$, $L_{r}, L_{c}, Ls_{r}, Ls_{c}, p$, over
\Ensure Collection $Ls_{r}, Ls_{c}, OPs_{r}, OPs_{c}$ that implement permutation matrix with length $d$
\State minm, $best_{r}^{cost}, best_{c}^{cost}$ $\leftarrow$ MAX \Comment{Assign the maximum value supported by the operating system as the initial upper bound.}
\State $visi_{r}, visi_{c}$, select\_list $\leftarrow$ list[$\varnothing$] \Comment{Initialize all required lists.}
\State LIMIT $\leftarrow$ LATEST\_MINM\_DEPTH \Comment{Define the previous minimum cost.}
\State close\_permu $\leftarrow$ False
\State stuck\_counter $\leftarrow$ 0, max\_stuck\_iterations $\leftarrow$ threshold \Comment{Define the local minima threshold.}
\While{$M$ $\neq$ Permutation}
	\State select\_list.clear()
	\State $L_{r}$.clear(), $L_{c}$.clear(), $L_{r}^{cost}$.clear(), $L_{c}^{cost}$.clear(), $B_{r}$.clear(), $B_{c}$.clear() \Comment{Clear all lists for the next iteration.}
	
	\State $H_{r}, H_{c}$, minm $\leftarrow$ functionCost($M, M', p$) \Comment{Calculate the current cost of $M$ and $M'$.}
    
	\State $L_{r} \leftarrow$ L\_Collection($visi_{r}$) \Comment{Collect row operations.}
	\State $B_{r}, best_{r}^{cost} \leftarrow$ Ops\_Sel($L_{r}, L_{r}^{cost}, M, M', p$, minm) \Comment{Select row operators for later comparison.}
	
	\State $L_{c} \leftarrow$ L\_Collection($visi_{c}$) \Comment{Collect column operations.}
	\State $B_{c}, best_{c}^{cost} \leftarrow$ Ops\_Sel($L_{c}, L_{c}^{cost}, M, M', p$, minm) \Comment{Select column operators for later comparison.}
	
	\State is\_stuck = $|B_{r}|$ == 0 $\land$ $|B_{c}|$ == 0 \Comment{Set to one if neither $B_{r}$ nor $B_{c}$ has any available operators.}
	
	\If{is\_stuck}
		\State stuck\_counter++ \Comment{Increase the stuck counter.}
	\Else
		\State stuck\_counter $\leftarrow$ 0 \Comment{Clear the counter when $B_{r}$ or $B_{c}$ has an available operator.}
    \EndIf
    
    \If{(stuck\_counter $\geq$ max\_stuck) || (is\_stuck $\land \sum{visi_{r}}$ == 0 $\land \sum{visi_{c}}$ == 0)} \Comment{The stall counter exceeds the local minima threshold and no more operations are available.}
		\State Candidates.clear()
		\State Candidates $\leftarrow$ Rand\_Ops\_Sel($L_{r}, L_{c}, M, M', p$) \Comment{Randomly select operators that increase the cost value, without needing to consider ``minm``.}
		\State minimum\_candidate $\leftarrow$ Candidates.sort() \Comment{Sort the candidates list from minimum to maximum.}
		\State select\_list $\leftarrow$ minimum\_candidate \Comment{Extract the minimum candidate into ``select\_list``.}
		\State minm $\leftarrow$ minimum\_candidate.cost \Comment{Update the cost to that of the minimum candidate operator.}
		\State stuck\_counter $\leftarrow 0$ \Comment{Clear the local minima counter.}
	\Else \Comment{When we do not fall into a local minimum, we select an available operator from $B_{r}$ or $B_{c}$.}
		\If{$best_{r}^{cost} < best_{c}^{cost}$}
			\State select\_list $\leftarrow$ $B_{r}$ \Comment{Choose row operations for ``select\_list``.}
			\State minm $\leftarrow$ $best_{r}^{cost}$ \Comment{Update the cost using the row operator.}
		\ElsIf{$best_{r}^{cost} > best_{c}^{cost}$}
			\State select\_list $\leftarrow$ $B_{c}$ \Comment{Choose column operations for ``select\_list``.}
			\State minm $\leftarrow$ $best_{c}^{cost}$ \Comment{Update the cost using the column operator.}
		\Else
			\State select\_list $\leftarrow$ $B_{c}$ if $|B_{c}| > 0 $ else $B_{r}$ \Comment{Prefer the column operators when $B_{r}^{cost}$ equal to $B_{c}^{cost}$ if $B_{c}$ exist, otherwise, we select $B_{r}$'s operators.}
			\State minm $\leftarrow$ $best_{c}^{cost}$ if |select\_list| $> 0$ else $best_{r}^{cost}$ \Comment{Update ``minm`` using $best_{c}^{cost}$ when ``select\_list`` selects from $B_{c}$; otherwise, update it using $best_{r}^{cost}$.}
    	\EndIf
    \EndIf
    
	\State $L_{r}, L_{c}, Ls_{r}, Ls_{c}, M, M', OPs_{r}, OPs_{c}, visi_{r}, visi_{c}, d$, over $\leftarrow$ Ops\_Exec()

	\If{close\_permu}
		\State $p \leftarrow -1$ \Comment{Switch the cost function to $H_{prod}$ when $M$ is close to a permutation matrix.}
	\EndIf
    
    \If{$d$ > LIMIT}
	    \State over $\leftarrow True$ \Comment{Break out of the iteration when the depth of $M$ exceeds the previous minimum cost value.}
	    \State \Return
    \EndIf
\EndWhile
\State \Return $L_{r}, L_{c}, Ls_{r}, Ls_{c}, M, OPs_{r}, OPs_{c}, d$, over
\end{algorithmic}
\end{algorithm}

\subsection{Parallel Greedy Algorithm}
The second strategy, the Parallel Greedy Algorithm
(Algorithm~\ref{alg:parallel-greedy}), aggregates both row and column
operations into a unified ``select\_list``. Similar to the Row-or-Col
greedy algorithm, this strategy also considers $B_{r}$ and $B_{c}$ to
maintain a bidirectional exploration of possible operators before
``close\_permu`` is triggered. The primary distinction of the
Parallel strategy, however, lies in its late-stage optimization: whereas
the Row-or-Col strategy switches to the $H_{prod}$ cost function once the
matrix approaches a permutation state, the Parallel strategy instead
restricts the search exclusively to column operations at this stage,
following the same ``close\_permu`` condition. That is, once the
``close\_permu`` flag is activated, we adopt the spirit of the
Row/Column greedy algorithm by considering only the available column
operators, rather than switching the cost function as in the Row-or-Col
strategy.

As in the Row-or-Col approach, if one operation type is unavailable, the
algorithm proceeds with the existing candidates. Likewise, a stagnation
threshold is used to escape local minima: if no progress is made for
three consecutive iterations, the algorithm accepts a cost-increasing
move.

In summary, we combine both $B_{r}$ and $B_{c}$ to expand the
``select\_list`` before the matrix reaches a low-Hamming-weight
state, selecting the minimum cost value between $B_{r}^{cost}$ and
$B_{c}^{cost}$, with column-wise operations considered first in our
approach. Together, the Row-or-Col and Parallel strategies constitute the
greedy algorithm parameters, allowing the matrix-specific strategy that
yields the optimization circuit to be selected in advance of synthesizing the
target quantum circuit.



\begin{algorithm}[htbp]
\caption{Parallel Greedy}
\label{alg:parallel-greedy}
\scriptsize
\begin{algorithmic}[1]
\Require $M \in \mathbb{F}_{2}^{n \times n}$, $M' \in \mathbb{F}_{2}^{n \times n}$, $L_{r}, L_{c}, Ls_{r}, Ls_{c}, p$, over
\Ensure Collection $Ls_{r}, Ls_{c}, OPs_{r}, OPs_{c}$ that implement permutation matrix with length $d$
\State minm, $best_{r}^{cost}, best_{c}^{cost}$ $\leftarrow$ MAX \Comment{Assign the maximum value supported by the operating system as the initial upper bound.}
\State $visi_{r}, visi_{c}$, select\_list $\leftarrow$ list[$\varnothing$] \Comment{Initialize all required lists.}
\State LIMIT $\leftarrow$ LATEST\_MINM\_DEPTH \Comment{Define the previous minimum cost.}
\State close\_permu $\leftarrow$ False
\State stuck\_counter $\leftarrow$ 0, max\_stuck\_iterations $\leftarrow$ threshold \Comment{Define the local minima threshold.}
\While{$M$ $\neq$ Permutation}
	\State select\_list.clear()
	\State $L_{r}$.clear(), $L_{c}$.clear(), $L_{r}^{cost}$.clear(), $L_{c}^{cost}$.clear(), $B_{r}$.clear(), $B_{c}$.clear() \Comment{Clear all lists for the next iteration.}
	
	\State $H_{r}, H_{c}$, minm $\leftarrow$ functionCost($M, M', p$) \Comment{Calculate the current cost of $M$ and $M'$.}
    
    \If{!close\_permu}
		\State $L_{r} \leftarrow$ L\_Collection($visi_{r}$) \Comment{Collect row operations when $M$ is not close to a low-Hamming-weight state.}
		\State $B_{r}, best_{r}^{cost} \leftarrow$ Ops\_Sel($L_{r}$, $L_{r}^{cost}$, $M$, $M'$, $p$, minm)  \Comment{Select row operators for later combination.}
    \EndIf
    
	\State $L_{c} \leftarrow$ L\_Collection($visi_{c}$) \Comment{Collect column operations; if $M$ is close to a permutation matrix, consider only these operators.}
	\State $B_{c}, best_{c}^{cost} \leftarrow$ Ops\_Sel($L_{c}$, $L_{c}^{cost}$, $M$, $M'$, $p$, minm) \Comment{Select column operators for later combination.}
	
	\State is\_stuck = $|B_{r}|$ == 0 $\land$ $|B_{c}|$ == 0
	
	\If{is\_stuck}
		\State stuck\_counter++ \Comment{Increase the stuck counter.}
	\Else
		\State stuck\_counter $\leftarrow$ 0 \Comment{Clear the counter when $B_{r}$ or $B_{c}$ has an available operator.}
    \EndIf
    
    \If{(stuck\_counter $\geq$ max\_stuck\_iterations) || (is\_stuck $\land \sum{visi_{r}}$ == 0 $\land \sum{visi_{c}}$ == 0)}
	    \State Candidates.clear()
		\State Candidates $\leftarrow$ Rand\_Ops\_Sel($L_{r}, L_{c}, M, M', p$) \Comment{Randomly select operators that increase the cost value, without needing to consider ``minm``.}
		\State minimum\_candidate $\leftarrow$ Candidates.sort() \Comment{Sort the candidates list from minimum to maximum.}
		\State select\_list $\leftarrow$ minimum\_candidate \Comment{Extract the minimum candidate into ``select\_list``.}
		\State minm $\leftarrow$ minimum\_candidate.cost \Comment{Update the cost to that of the minimum candidate operator.}
		\State stuck\_counter $\leftarrow 0$ \Comment{Clear the local minima counter.}
	\Else \Comment{When we do not fall into a local minimum, we select an available operator from $B_{r}$ + $B_{c}$.}
		\State select\_list $\leftarrow$ $B_{r} + B_{c}$
		\If{$best_{r}^{cost} < best_{c}^{cost}$}
			\State minm $\leftarrow$ $best_{r}^{cost}$
		\ElsIf{$best_{r}^{cost} > best_{c}^{cost}$}
			\State minm $\leftarrow$ $best_{c}^{cost}$
		\Else
			\State minm $\leftarrow$ $best_{c}^{cost}$ if |select\_list| $> 0$ else $best_{r}^{cost}$
    	\EndIf
	\EndIf
    
	\State $L_{r}, L_{c}, Ls_{r}, Ls_{c}, M, M', OPs_{r}, OPs_{c}, visi_{r}, visi_{c}, d$, over $\leftarrow$ Ops\_Exec()
	
	\If{close\_permu}
		\State $p \leftarrow -1$ \Comment{Switch the cost function to $H_{prod}$ when $M$ is close to a permutation matrix.}
	\EndIf
    
    \If{$d$ > LIMIT}
	    \State over $\leftarrow True$ \Comment{Break out of the iteration when the depth of $M$ exceeds the previous minimum cost value.}
	    \State \Return
    \EndIf
\EndWhile
\State \Return $L_{r}, L_{c}, Ls_{r}, Ls_{c}, M, OPs_{r}, OPs_{c}, d$, over
\end{algorithmic}
\end{algorithm}
